\documentstyle[german]{article}
\begin{document}

\pagestyle{myheadings} \markright{Rewrite} \pagenumbering{roman}
\title{Rewrite --- ein Mathematiksystem}
\author{
Rainer Heintzmann\thanks{Blumenhaller Weg 138,49078 Osnabr"uck}
}
\maketitle
\begin{abstract}
Rewrite ist ein Sytem zur Beieinflu"sung symbolischer Ausdr"ucke,
und zum Durchf"uhren numerischer Rechnungen.
Mittels Umschreiben der Ausdr"ucke in einem Regelsystem von `Algebren'
lassen sich einfache algebraische Funktionen wie Differenzieren, L"osen
von quadratischen Gleichungen usw. von Termen ausf"uhren.
Die Regelsysteme sind selbst zu definieren. Unter dem Verzeichnis `inc' 
sind einige Beispiele.

Numerisches Auswerten der Terme ist auch m"oglich. Benutzer k"onnen 
eigene Funktionen schreiben und sie in die Liste vorhandener
Funktionen in dem Modul `userfkt.c' einf"ugen.
\end{abstract}

\section{Starten des Programms}
Starten sie das Programm `rewrite' von der shell aus.
Geben Sie nun `INC ''inc/all''' ein. Es werden nun eine Reihe von Include-Files
geladen. Nun k"onnen Sie Ihre Aufgabe stellen z.B:
`REWR D(sin cos x,x) : All'
Das Programm bestimmt nun die Ableitung der Funktion `sin(cos(x))' nach x durch 
Umschreiben dieses Terms in der Algebraliste `All'.

\section{Benutzung von `Rewrite'}

Generell wird die Eingabe immer Zeilenweise verarbeitet. Bei gr"o"seren
Bl"ocken mu"s die Zeile durch ein `$\backslash$' am Zeilenende fortgesetzt werden.
Wichtig ist, da"s kein Zeichen (auch kein Leerzeichen) hinter dem `$\backslash$' steht.
Kommentare beginnen mit `\#' und gehen bis zum Ende der Zeile oder bis zum
n"achsten `$\backslash$' gefolgt von einem Zeilentrenner.
\newpage
\section{Liste aller Rewritebefehle}
\begin{description}
\item[Befehl] Erkl"arung
\item[INC ''filename'']
								Unter `filename' abgespeicherte Programme werden ausgef"uhrt.
                `Filename' ist ein von `''' umschlossener String.
                Mehrfache INCLUDES sind m"oglich bis zu einer maximalen Tiefe. 

                Bsp.: INC `'inc/all'' : alle wichtigen Include-files werden 
                gelesen.

                Der INC-Befehl wird einmalig beim verarbeiten der Eingabe
                ausgef"uhrt. Zeilen, die `INC' enthalten d"urfen nicht mit
                `$\backslash$' fortgesetzt werden. (Einzige Ausnahme !)

\item[REWR term : algname [;algname ..\symbol{93}]
                Der Term `term' wird in den Algebren (Regellisten) `algname'
                nacheinander umgeschrieben.

\item[PRTNUM ''String'']
								Das Standartformat f"ur die Ausgabe von reelen Zahlen wird in
                `String' ge"andert. Z.B. : PRTNUM ''\% 2.15g ''

\item[OUT ''filename'']
								Die Ausgabe wird in das File `filename' umgelenkt; dieses wird vorher
                gel"oscht. Zus"atzlich ist noch `OUT ''stdout''' m"oglich. Dadurch
                findet die Ausgabe wieder auf stdout (normalerweise der Bildschirm)
								statt. Bei jeder Ausgabe"anderung wird die noch offene Datei
								geschlossen.
                Grafische Darstellung von Funktionen mittels anderer Programme
                ist so einfach m"oglich.

\item[DATA Const]
								Liest eine Real-Zahl aus der ge"offneten Datei in die Konstante
                `Const'. So lassen sich z.B. auch Me"swerte weiterverarbeiten.

\item[DATA ''filename'']
								Zus"atzlich zu der normalen Standarteingabe besteht die
                M"oglichkeit mit `DATA Const' aus einer Datei zu lesen.
                Diese Datei mu"s mit `DATA ''filename''' ge"offnet werden.
								W"ahrend des Lesens wird beim Erreichen von nicht numerischen Daten
								der Rest, bis zum Zeilenende "ubersprungen.

\item[SYS ''String'']
								f"urt den in `String' angegebenen Befehl in der shell aus.

\item[SORT name [,name ..\symbol{93}]
								Erzeugung eines Datentyps. Der Typ wird bei der Definition von
								Funktionen, Variablen und Konstanten gebraucht.

\item[NUM sortname]
								Standartm"a"sig wird mit dem Datentyp `double' gerechnet
                Alle eingegebenen Zahlen (z.B.: 1.234) werden als `double'
                abgespeichert. Sie m"ussen aber einem selbstdefiniertem
                Datentyp (mit `SORT typname' erzeugt) zugeordnet werden.

                Bsp.: SORT real

											NUM real    \# ordnet Nummern den Typ real zu.

\item[CONST name [,name ..\symbol{93} : sortname]
								Deklaration einer Konstanten. Es wird eigentlich
                eine Konstante erzeugt, die beliebige Werte annehmen kann.
                (Also eine Variable). Im Gegensatz zu mit VAR erzeugten
								Metha-Variablen kann sie aber nicht als Ersatz f"ur ganze Terme
                benutzt werden.

\item[VAR name [,name ..\symbol{93} : sortname]
								Erzeugung einer Metha-Variablen. Man kann diese Variable
                als Platzhalter in Spezifikationen f"ur Algebren (Regellisten)
								verwenden.

                Bsp.:
											VAR A : real  SPEC xyz A+0=A

\item[NUM name [,name ..\symbol{93} : sortname]  
								Erzeugung einer Metha-Variablen. Man kann diese Variable
                als Platzhalter nur f"ur Numerische Werte in Spezifikationen 
                f"ur Algebren (Regellisten) verwenden.

                Bsp.: NUM N : real 

											SPEC xyz N+0=N

\item[FUNC name : sortlist \frq \frq\ sortname [LA1..5|RU\symbol{93} [ATT name\symbol{93}]
                Erzeugung einer Funktion. Die erste `sortlist' gibt die
                Anzahl und Typen der Argumente an. Es mu"s mindestens
								eines sein. Sie werden jeweils durch Kommata voneinander getrennt.
								Ein ++ hinter dem letzten Argument steht f"ur variabel viele Argumente des 
								letzen Types. Der andere `sortname' gibt den
                Resulttattyp der Funktion an. LA1 bis LA5 legt das Ein- bzw.
                Ausgabeformat der zu definierenden Funktion fest. Es bedeuten: \\
                   {\bf LA1 bis LA5}   : Links-assoziativer Operator (2 Argumente)
                                   LA5 klammert vor LA4 usw. \\
                   {\bf RU}            : Rechts-un"arer Operotor. (z.B. Fakult"at) \\
									 {\bf RA}            : Rechts-assoziativer Operator (z.B. a\^b)\\

                Ist kein spezielles Ein/Ausgabeformat angegeben, so wird die
                Funktion in der Form : `sin x' , `sin(x)' bzw. `f(x,y)'
                eingelesen.

                {\bf ATT name} legt den Namen der Funktion in der Liste 
                `userfunction' im Modul userfkt.c fest, die aufgerufen wird, 
                wenn alle Argumente der Funktion einen numerischen Wert besitzen.
								`name' ist nicht in `''' zu setzen. 
								Name darf frei gew"ahlt werden (auch gleichlautend mit
								einem Funktionsnamen).

                z.B.: FUNC + : real,real \frq \frq\ real LA1 ATT 2

								auch m"oglich ist: \\
							  {\bf ATT ( Argdefs ) Block}
								H"angt eine selbstdefinierte Funktion an. ( Argdefs ) sind 
								vollst"andige Definitionen f"ur die Argumente ( Konstanten,Terme
								oder Funktionen ).
							  Der entsprechende Funktionsk"orper wird ausgef"uhrt, wenn eine
								Zuweisung eines Terms an eine Konstante erfolgt, oder ein Term
								ausgegeben werden soll.

							  z.B. FUNC test : real,real,real \frq \frq\ real ATT $\backslash$ \\
							       (CONST a,b : real; TERM c : real) \{ RETURN c+a*b \}

\item[RETURN Term]
								Setzt den Resultatwert (einen Term) der Funktion. Die Ausf"uhrung
                der Funktion wird NICHT!! beendet. Enth"alt dieser Term lokale
								Variablen, so gibt es eine Fehlermeldung, und die Funktion gild
								als nicht ausgef"uhrt.

\item[OTHER name : sortlist \frq \frq\ sortname [ATT name\symbol{93}]
                Definiert zu einem existierenden Funktionsnamen eine weitere
                Abbildung mit anderen Argument- und Resultattypen. 
								Funktionen k"onnen anhand ihrer Argumenttypen unterschieden werden.

                z.B.: OTHER + : vector,vector \frq \frq\ vector

\item[BRA  auf zu : sortlist \frq \frq\ sortname [ATT name\symbol{93}]
                Definiert eine Funktion, die in Klammerschreibweise benutzt
                werden kann. `auf' bezeichnet das Symbol f"ur Klammer auf,
                und `zu' das Symbol f"ur Klammer zu.
								Es ist folgende zusaetzliche Notation bei der Benutzung erlaubt:

								a[b,c] statt [a,b,c]           ([ = auf, ] = zu)
								damit lassen sich z.B. Vectoren mit a[5] indizieren.

                z.B: BRA \flq \frq : vector,vector \frq \frq\ real \# Skalarprodukt \flq a ,b \frq \\
                Kennzeichnend f"ur OTHER ist die erste Klammer.\\
                z.B.: OTHER \flq : compl,compl \frq \frq\ real\\

\item[ALG algname] Definiert eine neue Algebra (Regelliste) , die Spezifikationen enthalten kann

\item[SPEC algname : [Fkt ...\symbol{93} was = zuwas [DO Block\symbol{93}]
               Alte Spezifikationen werden gel"oscht, diese neue Regel wird eingetragen.

\item[SPEC algname : algnamelist]
               Alte Spezifikationen werden gel"oscht, `algname' kann anstelle der
               gesamten algnamelist verwendet werden. In der
							 `algnamelist' stehen `algnames', die durch Kommata voneinander
							 getrennt sind. Hintereinander Umschreiben
               in mehreren Algebren kann so elegant geschrieben werden. 
               Prinzipiell d"urfen Listen von Algebren immer mit `(' `)' Klammern
               umgeben werden, was bedeutet, da"s beim Umschreiben die Algebren
               innerhalb der Klammern solange ausgef"uhrt werden, bis kein weiteres
               Umschreiben mehr m"oglich ist.

\item[ADDSPEC algname : [IF ( Fkt )\symbol{93} was = zuwas [DO Block\symbol{93}]
               Hinzuf"ugen einer Ersetzungsregel (Spezifikation) zur Algebra
                `algname'. Es d"urfen Metha-Variablen auf beiden Seiten
               verwendet werden. Mit NUM definierte Metha-Variablen werden
               nur bei Zahlen ersetzt.

							 Fkt ist selbstdefiniert (siehe IF ) oder eine der Funktionen :

                     LB(term,term)            `Lower or between' vergleicht 
                                              zwei Terme. \\
                     ISINT(term)              stellt fest ob term ein int-wert
                                              ist. \\
                     IN(terma,termb)          stellt fest ob terma in termb
                                              enthalten ist. \\
                     NIN(terma,termb)         verneinung von IN. \\

               Abh"angig vom Resultat der Funktionen wird die Regel angewandt
               oder eben nicht. Das Hintereinanderschreiben mehrerer Funktionen
               bewirkt eine und-Verkn"upfung der Resultate.

               Bsp.: ADDSPEC real : LB(A,B) A+B=B+A\\
               Ist DO in Verkn"upfung mit einem Block (oder einem Befehl)
               angegeben, so wird vor dem Ersetzen der linken durch die rechten
               Seite der Block ausgef"uhrt. Mehrere Befehle m"ussen in dem Block
							 von `\{' und `\}' umgeben sein.

               Bsp.: ADDSPEC real : Tay(A,B,D)=q DO \{ q=0 \\
                                    FOR i=1 TO N STEP 1 DO $\backslash$ \\
                                    q=q+D\^i/i!*Dp(B,D) \}\\

\item[LIFT (Fktlist(Args)) TO (Fktlist(Args))]
							 "Ubertragung s"amtlicher auf eine Menge von Funktionen und
							 Konstanten zutreffender Regeln auf eine andere solche Menge.
							 Fktlist(Args) ist eine durch Kommas getrennte Liste von Funktionen
							 mit irgentwelchen Argumenten des gew"unschten Datentyps oder
							 Konstanten.

               Die jeweiligen Funktionen werden durch den Typ ihrer Argumente
               eindeutig charakterisiert. Die Argumente spielen sonst im weiteren
               keine Rolle. Als Konstanten sind auch die reellen Zahlen 1 und 0
               m"oglich.

							 Alle Regeln, die in der linken Menge gelten werden
							 nun automatisch auch auf die rechte Menge angewandt.

							 Bsp. :\\
\begin{verbatim}
LIFT (Aag+Aag,Aag-Aag,-Aag,0ag) TO (Ar+Ar,Ar-Ar,-Ar,0r)
LIFT (Aag+Aag,Aag-Aag,-Aag,0ag) TO (Ar*Ar,Ar/Ar,mi(Ar),1r)
\end{verbatim}
               mit diesem Trick lassen sich die Eigenschaften eines
               Ringes leicht aufgrund schon spezifizierter Gruppeneigenschaften
               definieren.

\item[TERM name : sortname]
               Deklaration einer Term-Variable.
               Eine Term-Variable kann einen beliebigen Ausdruck des gleichen
               Typs aufnehmen.Ineianderschachtelung von Term-Variablen bewirkt
               ein Einsetzen der aktuellen Ausdr"ucke. Die Abh"angigkeit zu der anderen
               Term-Variablen wird nicht gespeichert.

               Immer vorhanden ist der Term `\%'. Er enth"alt den jeweils
               zuletzt ausgegebenen Ausdruck. Soll dieser weiterhin gespeichert
               werden, so mu"s er an eine andere Term-Variable zugewiesen werden.

               Term-Variablen werden eigentlich im Sinne von Funktionen
               verwendet. So ist die Zuweisung `TERM q  q=x\^y*a' im Sinne von 
               `q(x,y,a)' zu verstehen. Durch Zuweisungen an die einzelnen
               Variablen `x=4 y=2 a=3' kann dann einfach der Funktionswert von
               q(4,2,3) mit `q' ausgerechnet werden.
               

\item[ATT sortname Name1 , Name2 , Name3 , Name4] 
               Einem Typ wird Anzeige- ,Vergleichs- , Kopier- und L"oschfunktion
               f"ur numerische Werte zugewiesen. Dei `Namen' stehen hierbei nicht
							 in doppelten Anf"uhrunsstrichen.

               `Name1' bezeichnet die Anzeigefunktion in `userfunction' mit der
               numerische Werte vom Typ sortname angezeigt werden sollen. \\
               `Name2' bezeichnet die Vergleichsfunktion mit der zwei numerische
               Werte verglichen werden (1 : Zahl1 \frq\ Zahl2,0: Zahl1 == Zahl2,
               -1 : Zahl1 \flq Zahl2). \\
               `Name3' bezeichnet die Kopierfunktion in `userfunction' mit der
               Numerische Werte kopiert werden. \\
               `Name4' bezeichnet die L"oschfunktion in `userfunction' mit der
               Numerische Werte gel"oscht werden.

               Bsp.: ATT complex show\_complex,cmp\_complex,cpy\_complex,del\_complex\\
               siehe Modul userfkt.c


\item[FOR Constname = N1 TO N2 [STEP N3\symbol{93} DO Block]
							For-Schleife.

							Bsp.: FOR i = 20 TO 10 STEP -1 DO i

\item[IF ( Fkt ) Block [ ELSE Block \symbol{93}] 
              IF-Bedingung. Fkt ist der Name einer Funktion mit double Resultat, die
              TRUE/FALSE (1.0/0.0) zur"uckgibt.

              Bsp.: CONST inc\_xxx : real \\
                    inc\_xxx = 1\\
                    IF (inc\_xxx \frq 0) INC ''inc/xxx''\\

\item[WHILE ( Fkt ) Block]
              WHILE-Schleife. F"ur `Fkt' siehe `IF'.

              Bsp.: i = 1\\
                    WHILE (10\frq i) 2\verb|^|i

\item[Zuweisungen: Termname = term]
               term wird in Termname gespeichert

               Bsp.: TERM q : real   \\
							 q=1+1+a

               Der letzte ausgegebene Ausdruck ist unter `\%' als Termname
               immer verf"ugbar. Er wird also automatisch an `\%' zugewiesen.
               In selbstdefinierten Funktionen d"urfen lokale Variablen nicht an globalere Terme zugewiesen werden.
               (daraus resultiert eine Fehlermeldung bei der Ausgabe. (Zuweisung an \%))
               Werden Term-Variablen an Term-Variablen zugewiesen, so wird der
               Inhalt kopiert und gegebenenfalls numerisch ausgewertet. Wird also
               eine Term-Variable ver"andert hat dies keinen einflu"s auf andere
               Term-Variablen.


\item[Constname = value]
               Die Konstante `Constname' bekommt einen Wert

               z.B.: x=1+1+3    (wenn + mit Funktion verbunden )

               Steht bei `value' kein Wert (z.B. andere Konstante ohne
               Wert oder Funktion, die nicht rechnet) so wird der Wert der 
               Konstanten gel"oscht und sie erscheint in der Ausgabe wieder
							 mit ihrem Namen `Constname'.

               Sowohl eine Const- als auch eine Term-Variable k"onnen durch einfaches
               hinschreiben des Namens ausgegeben werden.
							 Ausgegebene Ausdr"ucke werden an `\%' zugewiesen.

							 zur Verdeutlichung ein Beispiel :



\begin{verbatim}
1+1 
2                  # wird ausgegeben 
TERM q : real 
q=1+1 
q 
2                  # wird ausgegeben (q bleibt aber erhalten) 
Der letzte Ausgedruckte Wert ist unter dem Term \% wieder zu erreichen.
Bestehende Abh"angigkeiten zu anderen Termen werden hier ignoriert.
Bsp.:
TERM q
q=b
a+q
a+b                # wird ausgegeben
q=c
%+2
a+b+2              # wird ausgegeben
\end{verbatim}
\end{description}

\section{Allgemeines}
Werden von `''' begrenzte Strings hingeschrieben, so wird der Text ausgegeben.
Sie k"onnen nicht zugewiesen werden. (ein String darf auch `''' enthalten.)

Eine Reihe `benutzerdefinierter' Funktionen sind schon implementiert.
Dazu geh"oren wichtige Funktionen auf reelen Zahlen (real.c),
complexen Zahlen (complex.c), reelen Vectoren (rvector.c) und reelen
Matrizen (rmatrix.c) und Verbindungen der Datentypen. Alle Funktionen
sind im Modul (userfkt.c) eingetragen.

Neue Funktionen m"ussen immer
ein array von Zeigern auf Argumente als Argument bekommen und einen Zeiger
auf das Resultat zur"uckliefern. Beispiele sind in den genannten c-Modulen
zu finden.

Die grundlegenden mathematischen Datentypen, Funktionen und Regeln sind
im Modul `alg' im Verzeichni"s `inc' definiert. Dort kann man auch nachlesen
welche mathematischen Funktionen definiert sind.\\ \\
{\bf Wichtigste Fehler des Systems:} \\
Der Speicher wird nur unzureichend freigegeben.\\
{\bf M"angel des Systems:} \\
W"unschend w"ahre ein mehr Objekt-orientierter Ansatz; einfache Dinge wie
differenzieren usw. sind dann deutlich schneller (verschicken von Nachrichten).
Es existieren keine bisher `Container-Datentypen' wie Arrays, Matritzen, usw. .
Diese m"ussen durch entsprechende Funktionen nachgebildet werden, was Rechenzeit
kostet.

\begin{appendix}
\section{Zusammenfassung der yacc-syntax}

\begin{verbatim}

%token			CST VAR SORT FUNC ALG SPEC ADDSPEC REWR PRT GRAPH PFEIL TERM NUM ATT IS EQ
%token			LA1 LA2 LA3 LA4 LA5 LOCAL RU RA KOMMU LB IN NIN PARAM FOR TO STEP DO DBG LAST
%token			BRA OTHER DATA OUT INC EXE VARPAR SYS IF ELSE WHILE PRTNUM LIFT QUIT RETURN
%token	<y\_sym>	Id String
%token	<y\_sym>	Sort Alg Term
%token	<y\_sym>	La1 La2 La3 La4 La5 Ru Ra Func Const Var Bra Bra1 Bra2
%token			Terms Form Sorts Functs Forms Algs Cond2 Cond
%type	<y\_int>	VORRANG
%type	<y\_nod>	vinstr definition instruction aninstruction instructions prntexpr  tinstructions oinstructions
%type	<y\_nod> idlist l\_thing l\_list sortlist vsortlist funclist formlist termlist alglist Alg2 prntlist
%type	<y\_nod>	term func step att namedId d\_list
%type	<y\_nod> formula cformula condition condfunc

%left	La1 
%left	La2 
%left	La3 
%left	La4 
%left	La5
%right	Ra 
%right	Ru 
%left	Func 
%%


prog		: line
				| prog line

line		: instructions '\n' 
													
definition	: FUNC namedId ':' vsortlist PFEIL Sort	VORRANG att 
			| OTHER func ':' vsortlist PFEIL Sort att 
			| ATT Sort Id ',' Id ',' Id ',' Id 
			| BRA Id Id	':'  vsortlist PFEIL Sort	att 
			| VAR idlist ':' Sort			 
			| NUM Sort	 			  			 
			| NUM idlist ':' Sort			 
			| CST idlist ':' Sort			 
			| SPEC Alg ':' alglist	   
			| SPEC Alg ':' formlist		 
			| ADDSPEC Alg ':' formlist 
			| SORT idlist					     

			| ALG idlist	             
			| TERM idlist ':' Sort	   

idlist	: namedId						
			| namedId ',' idlist		

l_thing : Const					
			| Term						

l_list	: l_thing			     
			| l_thing ',' l_list 

att 		: /* NULL */		
			  | ATT namedId		
				| ATT '(' 
								 d_list ')' instruction 

namedId : Id           
        | Const        
        | Var	         
        | func         
        | Sort         
        | Alg          
        | Term         


VORRANG		: /* nothing */				
			| LA1						
			| LA2						
			| LA3						
			| LA4						
			| LA5						
			| RA						
			| RU						

condition : condfunc					

condfunc : LB '(' termlist ')'
			| IN '(' termlist ')' 	
			| NIN '(' termlist ')'	
			| '(' term ')'   				

instructions  : instruction	
							| instruction ';' instructions 

tinstructions : oinstructions	
							| oinstructions '\n' tinstructions 			

oinstructions : instructions			

trenner : '\n' 
				| ';'

d_list : definition						 
			 | definition ';' d_list 

instruction : aninstruction     
			| definition		    			
			| '{'
				tinstructions '}' 			
			| WHILE condition instruction	
			| IF condition instruction	
			| IF condition instruction ELSE instruction	
			| FOR term EQ term TO term step DO instruction

aninstruction	: /* Nichts */ 	
			| SYS String					
			| INC String					
			| REWR term ':' alglist	
			| prntlist						
			| PRTNUM String				
			| PRT prntlist				
			| OUT String					
			| DATA String					
			| DATA Const					
			| Term EQ term				
			| Term EQ REWR term ':' alglist	
			| Const EQ term			
			| EXE func Id 			
			| LIFT '(' termlist ')' TO '(' termlist ')'	
			| DBG Const 	      	
			| RETURN term								
			| QUIT                      

step  : /* NULL */	
			| STEP term	

prntexpr : term						
			| String							

prntlist : prntexpr						
			| prntexpr ',' prntlist	

alglist		: Alg2						
			| Alg2 ',' alglist		

Alg2 : Alg									
			| '(' alglist ')'			

sortlist	: Sort							
			| Sort ',' sortlist			

vsortlist : sortlist					
			| sortlist VARPAR				

func		: Func				
			| La1						
			| La2						
			| La3						
			| La4						
			| La5						
			| Ra						
			| Ru						
			| Bra						
			| Bra1					

funclist	: func						
			| func ',' funclist		

term		: Term							
			| LAST								
			| Const								
			| Var				      		
			| '(' term ')'		    
			| term Ru				      
			| Func term				    
			| La1 term				    
			| La2 term				    
			| La3 term				    
			| La4 term				    
			| La5 term				    
			| Bra termlist Bra	  
			| Bra1 termlist Bra2	
			| Func '(' termlist ')'		
			| term La1 term		    
			| term La2 term				
			| term La3 term				
			| term La4 term				
			| term La5 term				
			| term Ra  term				

termlist	: term						
			| term ',' termlist		

formula		: term EQ term		

vinstr : /* */							
			| DO instruction			

cformula	: formula	vinstr					
          | IF condition formula	vinstr 

formlist	: cformula					
			/* | cformula ';' formlist 

%%
\end{verbatim}
\end{appendix}
\end{document}
